\documentclass[aspectratio=169,11pt]{beamer}
\usepackage{../phanes-beamer}
\renewcommand{\ModuleID}{08}
\title{Agentic engineering capstone}
\subtitle{Build, evaluate, and defend a bounded nuclear workflow}
\author{PHANES Initiative}
\institute{The University of Texas at Austin}
\date{September 12, 2026}
\begin{document}
\begin{frame}[plain]
\titlepage
\end{frame}
\begin{frame}[t,fragile]{Choose a workflow with observable success}
\fontsize{11.5}{14}\selectfont
\begin{itemize}\setlength{\itemsep}{2mm}
\item A calculation checker with independently known results.
\item A public-document review or requirements-traceability assistant.
\item An educational simulator with explicit physical limits.
\item A bounded analysis workflow with reproducible tool outputs.
\end{itemize}
\end{frame}
\begin{frame}[t,fragile]{Write a question the project can answer}
\fontsize{11.5}{14}\selectfont
\begin{itemize}\setlength{\itemsep}{2mm}
\item Identify the user, task, inputs, and required output.
\item Define the deterministic work and model-mediated work.
\item State exclusions and human decisions.
\item Specify a measurable baseline before building the agent.
\end{itemize}
\end{frame}
\begin{frame}[t,fragile]{Reuse the course's engineering artifacts}
\fontsize{10.5}{12.5}\selectfont\setlength{\tabcolsep}{4pt}\renewcommand{\arraystretch}{1.15}
\begin{tabular}{>{\raggedright\arraybackslash}p{\dimexpr 0.64\linewidth-10.24pt\relax}>{\raggedright\arraybackslash}p{\dimexpr 0.36\linewidth-5.76pt\relax}}
\toprule
\textbf{Artifact} & \textbf{Introduced in} \\ \midrule
Information boundary statement & Module 02 \\[2pt]
Skill input/output and failure contract & Module 03 \\[2pt]
Provenance record & Module 04 \\[2pt]
Evidence dossier & Module 05 \\[2pt]
Serving and workload report & Module 06 \\[2pt]
Corpus and citation record, if needed & Module 07 \\[2pt]
\bottomrule\end{tabular}\par\vspace{2mm}
\end{frame}
\begin{frame}[t,fragile]{Make milestones evidence-based}
\fontsize{10.5}{12.5}\selectfont\setlength{\tabcolsep}{4pt}\renewcommand{\arraystretch}{1.15}
\begin{tabular}{>{\raggedright\arraybackslash}p{\dimexpr 0.15\linewidth-3.6pt\relax}>{\raggedright\arraybackslash}p{\dimexpr 0.3\linewidth-7.2pt\relax}>{\raggedright\arraybackslash}p{\dimexpr 0.55\linewidth-13.2pt\relax}}
\toprule
\textbf{Weeks} & \textbf{Milestone} & \textbf{Evidence} \\ \midrule
1--2 & Proposal & Question, boundary, baseline, test plan \\[2pt]
3--6 & Build & Tools, skills and end-to-end runs \\[2pt]
7--8 & Harden & Failure tests and draft dossier \\[2pt]
9--10 & Peer review & Reproducible findings and responses \\[2pt]
11--12 & Handoff & Final report, demo and clean rerun \\[2pt]
\bottomrule\end{tabular}\par\vspace{2mm}
\end{frame}
\begin{frame}[t,fragile]{Review a proposal before coding}
\fontsize{11.5}{14}\selectfont
\begin{itemize}\setlength{\itemsep}{2mm}
\item Proposal: "Build an AI assistant for reactor design."
\item Narrow it to one input, one output, and one bounded workflow.
\item Name an independent reference and a failure test.
\item Identify the data boundary and the responsible reviewer.
\end{itemize}
\end{frame}
\begin{frame}[t,fragile]{One worked approach: narrowing the proposal}
\fontsize{11.5}{14}\selectfont
\begin{itemize}\setlength{\itemsep}{2mm}
\item One input: the public educational pin-cell geometry, specified by its XML and material data.
\item One output: a validated heat-balance calculation with stored result and provenance.
\item One bounded workflow: the agent runs the documented educational calculation package and checks the outcome against the analytical reference.
\item Failure test: an invalid input (negative flow) produces a nonzero exit and no success file.
\item Boundary and reviewer: all data is public or synthetic; the data inventory lists the inference endpoint as the only external destination; the instructor is the responsible reviewer.
\end{itemize}
\end{frame}
\begin{frame}[t,fragile]{Freeze an evaluation set before tuning}
\fontsize{11.5}{14}\selectfont
\begin{itemize}\setlength{\itemsep}{2mm}
\item Separate development examples from final evaluation tasks.
\item Include normal, boundary, invalid, and ambiguous cases.
\item Record metrics, acceptance criteria, and review rules.
\item Version the set when scope changes.
\end{itemize}
\end{frame}
\begin{frame}[t,fragile]{Run design review as an engineering exercise}
\fontsize{11.5}{14}\selectfont
\begin{itemize}\setlength{\itemsep}{2mm}
\item A peer team attempts reproducible technical and workflow failures.
\item Use public/synthetic cases and the agreed test environment.
\item Record the expected behavior, actual behavior, and consequence.
\item Tie severity to the project's declared use and boundary.
\end{itemize}
\end{frame}
\begin{frame}[t,fragile]{Four failure axes guide peer design review}
\fontsize{10.5}{12.5}\selectfont\setlength{\tabcolsep}{4pt}\renewcommand{\arraystretch}{1.15}
\begin{tabular}{>{\raggedright\arraybackslash}p{\dimexpr 0.22\linewidth-5.28pt\relax}>{\raggedright\arraybackslash}p{\dimexpr 0.44\linewidth-10.56pt\relax}>{\raggedright\arraybackslash}p{\dimexpr 0.34\linewidth-8.16pt\relax}}
\toprule
\textbf{Failure axis} & \textbf{Specific defect to test} & \textbf{Required review artifact} \\ \midrule
Technical & Unit mismatches, non-convergence, or geometry overlap. & Solver logs, statepoint hash and residual plots. \\[2pt]
Agentic & Hallucinated tool calls, runaway loops, or unverified output. & Harness execution trace and command history. \\[2pt]
Security & Prompt injection via inputs, broad permissions, or leaked keys. & OpenCode configuration and tool arguments. \\[2pt]
Information handling & Unapproved external egress, missing hashes, or boundary drift. & Data inventory, flow map and signed decision memo. \\[2pt]
\bottomrule\end{tabular}\par\vspace{2mm}
\end{frame}
\begin{frame}[t,fragile]{A useful finding is reproducible}
\fontsize{10.5}{12.5}\selectfont\setlength{\tabcolsep}{4pt}\renewcommand{\arraystretch}{1.15}
\begin{tabular}{>{\raggedright\arraybackslash}p{\dimexpr 0.27\linewidth-4.32pt\relax}>{\raggedright\arraybackslash}p{\dimexpr 0.73\linewidth-11.68pt\relax}}
\toprule
\textbf{Field} & \textbf{Required content} \\ \midrule
Identifier and severity & Unique label with impact rationale \\[2pt]
Preconditions & Exact version, data and configuration \\[2pt]
Steps & Commands or inputs to reproduce \\[2pt]
Expected / actual & Observable difference and artifacts \\[2pt]
Disposition & Fix, scope reduction, or documented acceptance \\[2pt]
\bottomrule\end{tabular}\par\vspace{2mm}
\end{frame}
\begin{frame}[t,fragile]{Respond to a stale-result finding}
\fontsize{11.5}{14}\selectfont
\begin{itemize}\setlength{\itemsep}{2mm}
\item The report reused a successful output after the input changed.
\item Identify the immediate containment action.
\item Specify a durable fix and a regression test.
\item Which earlier results need to be rerun?
\end{itemize}
\end{frame}
\begin{frame}[t,fragile]{One worked approach: containing a stale result}
\fontsize{11.5}{14}\selectfont
\begin{itemize}\setlength{\itemsep}{2mm}
\item Containment: mark the affected report and cases as invalidated, and stop citing those outputs in any downstream document.
\item Durable fix: bind result reuse to the input hash and solver identity, so a changed input cannot retrieve the older success record.
\item Regression test: the changed-input reuse test; a changed input hash must not return the previous result.
\item Rerun scope: rerun every result whose provenance may have mismatched, and record the rerun in the manifest and the dossier review response.
\end{itemize}
\end{frame}
\begin{frame}[t,fragile]{Make the handoff work on a clean machine}
\fontsize{11.5}{14}\selectfont
\begin{itemize}\setlength{\itemsep}{2mm}
\item Document Linux prerequisites and pinned project dependencies.
\item Provide .env.example with placeholders and startup instructions.
\item Include a public sample, expected result, and failure demonstration.
\item Retain versions, licenses, limitations, and ownership.
\end{itemize}
\end{frame}
\begin{frame}[t,fragile]{Grade evidence and engineering judgment}
\fontsize{10.5}{12.5}\selectfont\setlength{\tabcolsep}{4pt}\renewcommand{\arraystretch}{1.15}
\begin{tabular}{>{\raggedright\arraybackslash}p{\dimexpr 0.25\linewidth-6.0pt\relax}>{\raggedright\arraybackslash}p{\dimexpr 0.35\linewidth-8.4pt\relax}>{\raggedright\arraybackslash}p{\dimexpr 0.4\linewidth-9.6pt\relax}}
\toprule
\textbf{Component} & \textbf{Draft weight} & \textbf{What earns credit} \\ \midrule
Proposal & 10\% & Feasible scope and boundary \\[2pt]
Build milestones & 30\% & Working tools and integration \\[2pt]
Evaluation + hardening & 30\% & Relevant tests and honest limits \\[2pt]
Review response & 10\% & Reproduced findings and closure \\[2pt]
Final report + demo & 20\% & Traceable results and clean handoff \\[2pt]
\bottomrule\end{tabular}\par\vspace{2mm}
\end{frame}
\begin{frame}[t,fragile]{Final demonstration: show the complete chain}
\fontsize{11.5}{14}\selectfont
\begin{itemize}\setlength{\itemsep}{2mm}
\item State the engineering question and approved inputs.
\item Run the workflow from a clean, documented starting point.
\item Trace one result to its calculation or source evidence.
\item Demonstrate a failure, then defend the accepted scope.
\end{itemize}
\end{frame}
\begin{frame}[t,fragile]{Example capstone acceptance contract}
\fontsize{10.5}{12.5}\selectfont\setlength{\tabcolsep}{4pt}\renewcommand{\arraystretch}{1.15}
\begin{tabular}{>{\raggedright\arraybackslash}p{\dimexpr 0.26\linewidth-4.16pt\relax}>{\raggedright\arraybackslash}p{\dimexpr 0.74\linewidth-11.84pt\relax}}
\toprule
\textbf{Deliverable} & \textbf{Concrete acceptance evidence} \\ \midrule
OpenMC comparison & Two identity-bound runs, geometry/convergence review and uncertainty-aware comparison. \\[2pt]
Agent workflow & Scoped handoffs, permission demonstration and one verifier-found defect. \\[2pt]
Reproduction & A fresh environment can regenerate inputs and rerun an approved smoke case. \\[2pt]
Handoff & No embedded credentials; clear data prerequisites, limitations and owner. \\[2pt]
\bottomrule\end{tabular}\par\vspace{2mm}
\end{frame}
\begin{frame}[t,fragile]{What you should carry into engineering practice}
\fontsize{11.5}{14}\selectfont
\begin{itemize}\setlength{\itemsep}{2mm}
\item Use models where interpretation and orchestration help.
\item Keep calculations, source evidence, and interfaces inspectable.
\item Treat information handling and verification as design inputs.
\item Accept a system only for the scope its evidence supports.
\end{itemize}
\end{frame}
\begin{frame}[t,fragile]{Capstone exit ticket: defending an agentic workflow}
\fontsize{11.5}{14}\selectfont
\begin{itemize}\setlength{\itemsep}{2mm}
\item How do you demonstrate that the model did not hallucinate calculation outputs?
\item How is the project's regulatory information boundary verified and enforced?
\item What evidence demonstrates that the agent operated within bounded autonomy?
\item Under what specific conditions must the workflow stop and require human review?
\end{itemize}
\end{frame}
\begin{frame}[t,fragile]{A complete answer: the capstone defense}
\fontsize{11.5}{14}\selectfont
\begin{itemize}\setlength{\itemsep}{2mm}
\item Calculation defense: retain solver command, exit code, input and output hashes, and compare reported values directly to the binary or JSON output artifact.
\item Boundary enforcement: maintain an exact byte inventory and network flow map; verify that nonpublic data never traverses unapproved endpoints or external tools.
\item Bounded autonomy: prove that file edits and bash execution were constrained by reviewed permissions, finite loop budgets, and an independent verifier gate.
\item Mandatory escalation: stop when inputs violate physical ranges, tests reveal unhandled exceptions, new external destinations are requested, or regulatory scope is ambiguous.
\end{itemize}
\end{frame}
\begin{frame}[t]{References 1}
\fontsize{9}{12}\selectfont\begin{itemize}\setlength{\itemsep}{3mm}
\item \textbf{curriculum}: User-supplied PHANES curriculum: mod/01 through mod/08/README.md.
\item \textbf{colors}: \href{https://umac.utexas.edu/brand-center/colors/}{umac.utexas.edu/brand-center/colors}
\end{itemize}\end{frame}
\end{document}